\chapter{Global Kinematics: The Cosmic Mass and the Friedmann Replacement}
\label{chap:global_kinematics}

Standard cosmology treats the total mass and critical density of the universe as arbitrary parameters inherited from a primordial expansion. Tension Field Theory (TFT) does not accept this unconstrained metric expansion. Because the Spacetime Gravitational Floor (SGF) is a strictly bounded, visco-elastic mechanical structure, the mass and density of the observable universe are not arbitrary; they are strictly governed by the maximum gravitational strain the lattice can physically sustain.

\section{The Cosmic Mass Limit (Causal Saturation)}
As established in Part II, the dimensionless gravitational strain ($\alpha_G$) of the lattice is defined by the ratio of gravitational potential to the elastic limit of the vacuum ($c_E^2$). Therefore, to find the total mass of the observable universe ($M_U$), we consider the Type II Causal Saturation for total lattice strain.

While a Black Hole represents a Type I Causal Failure (a local density rupture), the observable universe represents a Type II Causal Saturation. The average density remains low, but the integrated gravitational potential ($\phi_{G,U}$) over the macroscopic cosmic radius ($R_U$) completely exhausts the lattice's cumulative energy budget:
\begin{equation}
    \phi_{G,U} = \int_0^{R_U} \frac{G \cdot dm}{r} = c_E^2
\end{equation}

At this absolute boundary limit of the Viscous Cosmic Horizon, the macroscopic strain reaches unity ($\alpha_G = 1$). Equating the effective global potential of the causal mass ($M_U$) at the optical drag limit ($R_U = c_E / \Gamma_D$):
\begin{equation}
    1 = \frac{GM_U}{R_U c_E^2} = \frac{GM_U}{\left(\frac{c_E}{\Gamma_D}\right) c_E^2}
\end{equation}

Solving for the total Cosmic Critical Mass ($M_U$):
\begin{equation}
    M_U = \frac{c_E^3}{G \Gamma_D}
\end{equation}

Substituting the established kinematic baseline for the Vacuum Damping Rate ($\Gamma_D \approx 2.3658 \times 10^{-18} \text{ s}^{-1}$) and standard physical constants, we evaluate the exact macroscopic mass limit:
\begin{equation}
    M_U = \frac{(299792458 \text{ m/s})^3}{(6.6743 \times 10^{-11} \text{ m}^3/\text{kg}\cdot\text{s}^2) \cdot (2.3658 \times 10^{-18} \text{ s}^{-1})} \approx 1.7064 \times 10^{53} \text{ kg}
\end{equation}

\textbf{The Universal Unification:} This cosmic mass limit reveals a structural symmetry within the TFT framework. In Part II, we derived the surface saturation Mass ($M_{\text{Sat}}$) of a Black Hole, defined as the point where the surface density dilutes to match the ambient vacuum and the Event Horizon dissolves. The mass required to saturate and dissolve the Black Hole horizon ($M_{\text{Sat}}$) is mechanically and mathematically identical to the critical mass bounding the observable universe ($M_U$):
\begin{equation}
    M_U = M_{\text{Sat}} = \frac{c_E^3}{G \Gamma_D} \approx 1.7064 \times 10^{53} \text{ kg}
\end{equation}

This mathematical equivalence proves that the observable universe is not expanding into an infinite void; its causal boundary ($R_U$) operates at the exact mechanical limit of a saturated Type I Black Hole. Consequently, the local causal domain is perfectly mechanically closed and structurally impervious to a kinematic ``Big Rip'' (runaway metric expansion). However, as established in the following chapters, the ultimate fate of this bounded system is strictly governed by the topological capacity of the vastly larger macroscopic reservoir ($M_{\text{Res}}$), culminating in a global Catastrophic Decompression (The Big Rupture) when the accumulation of frozen mass triggers lattice percolation.

\section{The Kinematic State Equation (The SGF Friedmann Replacement)}
Standard cosmology utilizes the Friedmann equation to link the expansion rate ($H_0$) to the critical mass density of the universe ($\rho_c$). TFT derives a mathematical equivalent formulation strictly from the solid-state geometry of this saturated strain limit.

Assuming the macroscopic mass is distributed uniformly, the critical mass density ($\rho_c$) of the observable universe is the total mass ($M_U$) divided by the geometric volume of the Viscous Cosmic Horizon ($V_U = \kappa R_U^3$), where $\kappa$ is the geometric projection parameter:
\begin{equation}
    \rho_c = \frac{\left( \frac{c_E^3}{G \Gamma_D} \right)}{\kappa \left( \frac{c_E}{\Gamma_D} \right)^3} = \frac{\left( \frac{c_E^3}{G \Gamma_D} \right)}{\kappa \frac{c_E^3}{\Gamma_D^3}}
\end{equation}

Canceling $c_E^3$ and algebraically isolating the kinematic damping ($\Gamma_D^2$) yields the TFT Kinematic State Equation:
\begin{equation}
    \Gamma_D^2 = \kappa G \rho_c
\end{equation}

This formulation naturally recovers the architecture of the standard Friedmann equation without invoking a stretching metric. The variance in the final coefficient is governed entirely by the geometric parameter ($\kappa$) derived in Part II.

For Newtonian radial projection within a 3D visco-elastic solid, $\kappa = 4\pi/3$. For relativistic tensor coupling compatible with General Relativity, $\kappa = 8\pi/3$. By applying the native 3D flat-space projection ($\kappa = 4\pi/3$), the TFT formulation becomes $\Gamma_D^2 = \frac{4\pi G}{3} \rho_c$.

Consequently, TFT's critical density is mathematically exactly double that of the standard model ($\rho_{c(\text{TFT})} = 2\rho_{c(\text{Friedmann})}$), reflecting the higher mass requirement to physically strain a high-tension solid compared to curving an empty geometric void.

\section{The Total Cosmic Reservoir ($M_{\text{Res}}$)}
\label{sec:reservoir_mass}

The critical mass calculated above ($M_U \approx 1.7064 \times 10^{53} \text{ kg}$) represents strictly the observable domain---the local optical and causal limit defined by the Viscous Cosmic Horizon ($R_U \approx 1.2672 \times 10^{26} \text{ m}$). However, as established in the Black Hole phase transition mechanism, the contiguous FCC lattice is vastly larger.

To form a saturated macroscopic core, the crystallization front must theoretically harvest a volume of active vacuum defined by the strict mass-energy conservation of the phase transition. In \hyperref[sec:reservoir_radius]{Section \ref*{sec:reservoir_radius}}, we derived the exact volumetric scaling factor between the observable causal domain and the ultimate structural capacity of the field using the phase density ratio:
$$ \frac{\rho_S}{\rho_{QG}} \approx 1.1633 \times 10^{26} $$

The total absolute mass of the cosmic reservoir is geometrically defined by multiplying the critical causal mass ($M_U$) by this exact volumetric multiplier:
$$ M_{\text{Res}} = M_U \times \left( \frac{\rho_S}{\rho_{QG}} \right) $$
$$ M_{\text{Res}} \approx (1.7064 \times 10^{53} \text{ kg}) \cdot (1.1633 \times 10^{26}) \approx 1.9851 \times 10^{79} \text{ kg} $$

This mechanically confirms that the observable universe ($10^{53} \text{ kg}$) is merely a single causal tile within a vastly larger, contiguous substrate. The macroscopic reservoir ($M_{\text{Res}} \approx 1.9851 \times 10^{79} \text{ kg}$) defines the ultimate topological mass capacity of the field, setting the absolute limit for the global thermodynamic clock before structural percolation triggers Catastrophic Decompression.

\section{The Baryonic Mass Cycle and the 5\% Volumetric Limit}
\label{sec:octahedral_masslimit}
Standard cosmology observes the baryonic mass density of the universe to be roughly $5\%$ ($\Omega_b$), treating it as an arbitrary parameter. Within Tension Field Theory, the compositional breakdown of the total cosmic mass ($M_{\text{Res}}$) is strictly dictated by the volumetric and topological limits of the FCC lattice. The existence and collapse of macroscopic mass operate within a strictly bounded phase cycle:

\begin{itemize}
    \item \textbf{Phase I (Kinetic Partitioning and Condensation):} As derived in Part I, macroscopic mass originates as a localized vibrational defect. To maintain mechanical equilibrium, this kinetic excitation reduces the fundamental static tension of the lattice nodes ($\beta_\Lambda$) from the ideal geometric packing limit of $\approx 74\%$ down to $\approx 69\%$. The total interstitial condensate filling the remaining voids ($\approx 26\%$) acts as the background hydrodynamic matrix, observationally identified as Dark Matter ($\beta_{IC}$). However, when baryonic mass physically condenses into macroscopic structures, its isotropic (spherical) geometry confines it strictly to the largest perfectly spherical interstitial pockets available---the Octahedral voids.
    
    \item \textbf{Phase II (Volumetric Saturation / Transition Trigger):} As local mass condenses, the density scales from the active vacuum baseline ($\rho_Q$) toward the convergence limit ($\rho_{QG}$). Because gravitational condensation is fundamentally isotropic, it cannot deform into the irregular, anisotropic capillary channels that make up the remaining $20.7\%$ of the void space. Therefore, the active vacuum can safely harbor condensed baryonic mass only up to the exact geometric threshold of the perfectly inscribed Octahedral voids ($\beta_{\text{oct}}$):
    \begin{equation}
        \beta_{\text{oct}} = \frac{\pi(\sqrt{2}-1)^3}{3\sqrt{2}} \approx 0.0526
    \end{equation}
    Exceeding this $5.26\%$ limit geometrically exhausts the perfectly spherical condensate capacity. Rather than undergoing extreme anisotropic deformation into the remaining capillary network, the system takes the path of least mechanical resistance, triggering hyperstatic structural failure.
    
    \item \textbf{Phase III (The Isoenergetic Phase Collapse):} When the local mass density reaches the void saturation trigger ($\rho_{QG}$), the active FCC lattice structurally yields, undergoing an instantaneous collapse into the zero-void Simple Cubic substrate ($\rho_S$). The kinematic degrees of freedom are locked into static structural stress, forming a State II Black Hole core. Because this ultimate phase transition is mathematically isoenergetic ($E_{QG} = E_S$), the observable baryonic mass capacity of the universe ($\beta_b$) is strictly bounded by this invariant geometric void limit:
    \begin{equation}
        \beta_b \le \beta_{\text{oct}} \approx 0.0526
    \end{equation}
\end{itemize}

(The complete derivation of the void saturation limit, topological percolation, and the isoenergetic phase continuity is detailed in \hyperref[app:phase_transition]{Appendix \ref*{app:phase_transition}} and \hyperref[app:energy_and_chronology]{Appendix \ref*{app:energy_and_chronology}}).